\documentclass[11pt]{article}
\usepackage[margin=1in]{geometry}
\usepackage{amsmath, amssymb}
\usepackage{graphicx}
\usepackage{enumitem}
\usepackage{booktabs}
\usepackage{hyperref}

\title{Method of Characteristics Nozzle Design:\\A Complete Walkthrough}
\author{}
\date{}

\begin{document}
\maketitle

\section{The Big Picture (ELI5 Version)}

Imagine you're at the throat of the nozzle. The flow is barely supersonic ($M \approx 1 + \varepsilon$). As air flows along the curved wall near the throat, the wall ``turns away'' from the flow, creating expansion waves. These waves bounce around inside the nozzle. Your job is to:

\begin{enumerate}
    \item Figure out where those waves go (using characteristic lines).
    \item Figure out the flow properties at every intersection point.
    \item Use those wave intersections with the upper wall to define the nozzle contour.
\end{enumerate}

Your project specifies a Mach 6 exit, a 2D planar nozzle with a circular-arc throat whose radius of curvature equals the throat half-height, and explicitly states \textbf{not} to compute a minimum-length nozzle with sharp corners~[1].

\section{Fundamental Concepts}

\subsection{What Are Characteristics?}

In supersonic flow, information travels along specific lines called \textbf{characteristics} (or Mach lines). There are two families:

\begin{itemize}
    \item \textbf{$C^-$ characteristics} (right-running): lines that slope at angle $(\theta - \mu)$ from horizontal.
    \item \textbf{$C^+$ characteristics} (left-running): lines that slope at angle $(\theta + \mu)$ from horizontal.
\end{itemize}

where:
\begin{itemize}
    \item $\theta$ = local flow angle (measured from horizontal). (Your ``$\delta
```'.)
    \item $\mu$ = Mach angle $= \arcsin(1/M)$.
\end{itemize}

\subsection{What Are the Riemann Invariants?}

Along each characteristic family, a specific combination of flow angle and Prandtl--Meyer function is \textbf{constant}:

\begin{align}
    \text{Along } C^-: \quad K^- &= \theta + \nu = \text{constant} \\
    \text{Along } C^+: \quad K^+ &= \theta - \nu = \text{constant}
\end{align}

These are your Riemann invariants. If you know $K^-$ and $K^+$ at a point, you know everything about the flow there:
\begin{align}
    \theta &= \frac{K^- + K^+}{2} \\
    \nu &= \frac{K^- - K^+}{2}
\end{align}

From $\nu$, you get $M$ (by inverting the Prandtl--Meyer function), and from $M$ you get $\mu$, $p/p_0$, $T/T_0$, etc.

\section{Your Specific Setup}

From the project document~[1]:
\begin{itemize}
    \item Exit Mach number: $M_e = 6$
    \item Nozzle type: 2D planar (rectangular cross-section, no axisymmetric correction)
    \item Throat shape: circular arc with radius of curvature $= h^*$ (throat half-height)
    \item \textbf{Not} a minimum-length nozzle (no sharp corner at throat)
    \item Exit half-height: $h_n = 175$~mm (half of 350~mm)
\end{itemize}

The expansion through the throat region is approximated with $N$ equally-spaced Prandtl--Meyer waves, and you cannot start exactly at the throat since $\mu = 90^\circ$ at $M=1$~[1].

\subsection{Known Quantities}
\begin{itemize}
    \item $\nu_e$ = Prandtl--Meyer function evaluated at $M_e = 6$
    \item $\theta_{\max} = \nu_e / 2$ (maximum wall turning angle, at the inflection point)
    \item $h^*$ and $h_e$ related by the isentropic area ratio $A_e/A^* = h_e/h^*$
\end{itemize}

\section{Sign and Geometry Conventions}

Consider the \textbf{upper wall} of the nozzle (symmetric about the horizontal centerline).

\begin{itemize}
    \item The $x$-axis is horizontal (flow direction).
    \item The $y$-axis is vertical (upward positive).
    \item $\theta > 0$ means the flow is angled upward.
    \item At the upper wall (curving away upward), $\theta > 0$.
\end{itemize}

\textbf{Characteristic slopes in physical space:}
\begin{align}
    C^+ \text{ slope:} \quad \frac{dy}{dx} &= \tan(\theta + \mu) \quad \text{(left-running, goes upward)} \\
    C^- \text{ slope:} \quad \frac{dy}{dx} &= \tan(\theta - \mu) \quad \text{(right-running, goes downward)}
\end{align}

From the upper wall, a $C^-$ characteristic goes \emph{downward} toward the centerline (since $\mu > \theta$ for low supersonic Mach, the slope $\tan(\theta - \mu)$ is negative). This makes physical sense: waves from the upper wall propagate downward into the flow.

\section{Complete Step-by-Step Procedure}

\subsection{Step 1: Initial Points on the Throat Wall}

Your $N$ expansion waves emanate from the upper wall of the circular-arc throat. Each wave is a $C^-$ characteristic going from the wall toward the centerline.

The wall angles are equally spaced:
\begin{equation}
    \Delta\theta = \frac{\theta_{\max}}{N} = \frac{\nu_e}{2N}
\end{equation}

At the $i$-th point on the throat wall ($i = 1, 2, \ldots, N$):
\begin{equation}
    \theta_i = i \cdot \Delta\theta
\end{equation}

\subsubsection{Why $\nu_i = \theta_i$ in the Initial Simple-Wave Region}

In the expansion region right after the throat, this is a \textbf{simple-wave region}. The $C^+$ characteristics passing through this region originate from near the sonic line (centerline region where $M \approx 1$). At the sonic line, $\nu = 0$ and $\theta = 0$, so:
\begin{equation}
    K^+ = \theta - \nu = 0 - 0 = 0
\end{equation}

Since $K^+ = 0$ everywhere in the simple-wave region, at every point:
\begin{equation}
    \theta - \nu = 0 \quad \Longrightarrow \quad \boxed{\nu = \theta}
\end{equation}

Therefore, at wall point $i$:
\begin{align}
    \theta_i &= i\,\Delta\theta \\
    \nu_i &= i\,\Delta\theta \\
    K^-_i &= \theta_i + \nu_i = 2i\,\Delta\theta \\
    K^+_i &= 0
\end{align}

From $\nu_i$, compute $M_i$ (invert the Prandtl--Meyer function), then $\mu_i = \arcsin(1/M_i)$.

\subsection{Step 2: Physical Locations of Initial Wall Points}

The throat arc has radius $R_{\mathrm{arc}} = h^*$. Place the origin at the centerline at the throat. The arc center for the upper wall is at $(0,\; h^* + R_{\mathrm{arc}}) = (0,\; 2h^*)$.

At angular position $\alpha$ measured from the downward vertical (so $\alpha = 0$ is the throat minimum and the wall angle equals $\alpha$):
\begin{align}
    x_{\mathrm{wall}} &= h^* \sin(\alpha) \\
    y_{\mathrm{wall}} &= h^*(2 - \cos\alpha)
\end{align}

For starting point $i$ with wall angle $\theta_i = i\,\Delta\theta$, set $\alpha = \theta_i$:
\begin{align}
    x_i &= h^* \sin(i\,\Delta\theta) \\
    y_i &= h^*(2 - \cos(i\,\Delta\theta))
\end{align}

\subsection{Step 3: Reflections off the Centerline (Symmetry Line)}

When a $C^-$ characteristic hits the centerline, it ``reflects'' as a $C^+$ characteristic. By symmetry, $\theta = 0$ on the centerline.

When $C^-$ char \#$i$ (carrying $K^-_i = 2i\,\Delta\theta$) reaches the centerline:
\begin{align}
    K^-_i &= \theta + \nu = 0 + \nu \quad \Longrightarrow \quad \nu = 2i\,\Delta\theta \\
    K^+_{\mathrm{reflected}} &= \theta - \nu = 0 - 2i\,\Delta\theta = -2i\,\Delta\theta
\end{align}

So the reflected $C^+$ characteristic carries $K^+ = -2i\,\Delta\theta$.

\subsection{Step 4: Interior Points (Intersections of $C^-$ and $C^+$ Characteristics)}

When a $C^-$ carrying $K^-_a$ and a $C^+$ carrying $K^+_b$ intersect at point $P$:
\begin{align}
    \theta_P &= \frac{K^-_a + K^+_b}{2} \\[6pt]
    \nu_P &= \frac{K^-_a - K^+_b}{2}
\end{align}

From $\nu_P$: compute $M_P$ (invert Prandtl--Meyer). Then $\mu_P = \arcsin(1/M_P)$.

\subsection{Step 5: Physical Location of Interior Points}

Suppose point $P$ is the intersection of:
\begin{itemize}
    \item A $C^-$ characteristic coming from point $A$ at $(x_A, y_A)$
    \item A $C^+$ characteristic coming from point $B$ at $(x_B, y_B)$
\end{itemize}

You already know $\theta$ and $\mu$ at all three points ($A$, $B$, and $P$). Use \textbf{average slopes}:
\begin{align}
    \text{Slope of } C^- \text{ from } A \text{ to } P: \quad m^- &= \tan\!\left(\frac{\theta_A + \theta_P}{2} - \frac{\mu_A + \mu_P}{2}\right) \\[6pt]
    \text{Slope of } C^+ \text{ from } B \text{ to } P: \quad m^+ &= \tan\!\left(\frac{\theta_B + \theta_P}{2} + \frac{\mu_B + \mu_P}{2}\right)
\end{align}

Now solve two simultaneous linear equations:
\begin{align}
    y_P - y_A &= m^-(x_P - x_A) \\
    y_P - y_B &= m^+(x_P - x_B)
\end{align}

Solving for $x_P$:
\begin{equation}
    x_P = \frac{(y_B - y_A) + m^- x_A - m^+ x_B}{m^- - m^+}
\end{equation}
Then:
\begin{equation}
    y_P = y_A + m^-(x_P - x_A)
\end{equation}

\subsection{Step 6: Wall Points in the Straightening Section}

After the inflection point (where $\theta = \theta_{\max}$), the wall must straighten to cancel all remaining waves. When a $C^+$ characteristic arrives at the upper wall, it is \textbf{absorbed} (no reflection).

\subsubsection{The No-Reflection Condition}

For uniform exit flow, we need $\theta = 0$ and $\nu = \nu_e$ at the exit. This means the $C^-$ invariant throughout the straightening section's wall must be:
\begin{equation}
    \boxed{K^-_{\mathrm{wall}} = \nu_e \quad \text{(constant for all wall points in the straightening section)}}
\end{equation}

\textbf{Why?} If the exit flow is uniform with $\theta = 0$ and $\nu = \nu_e$, then $K^- = 0 + \nu_e = \nu_e$ everywhere in the uniform region. Since $K^-$ is constant along $C^-$ characteristics, any $C^-$ characteristic \emph{leaving} the straightening-section wall must carry $K^- = \nu_e$. This means no new expansion is generated at the wall.

\subsubsection{Computing Wall Point Properties}

At wall point $j$ in the straightening section, a $C^+$ characteristic arrives carrying $K^+_j$ (known from the interior computation). Then:
\begin{align}
    \theta_{\mathrm{wall},j} &= \frac{K^-_{\mathrm{wall}} + K^+_j}{2} = \frac{\nu_e + K^+_j}{2} \\[6pt]
    \nu_{\mathrm{wall},j} &= \frac{K^-_{\mathrm{wall}} - K^+_j}{2} = \frac{\nu_e - K^+_j}{2}
\end{align}

Since $K^+_j < 0$ (from Step 3), $\theta_{\mathrm{wall},j}$ is positive but \textbf{less than} $\theta_{\max}$, and it \textbf{decreases} toward zero as you move downstream. The wall straightens out.

\subsubsection{Finding the Physical Location of Wall Points}

You know:
\begin{itemize}
    \item The previous wall point at $(x_{\mathrm{prev}}, y_{\mathrm{prev}})$ with angle $\theta_{\mathrm{prev}}$
    \item The interior point $(x_{\mathrm{int}}, y_{\mathrm{int}})$ from which the $C^+$ arrives
\end{itemize}

Wall segment slope (average):
\begin{equation}
    m_{\mathrm{wall}} = \tan\!\left(\frac{\theta_{\mathrm{prev}} + \theta_{\mathrm{wall},j}}{2}\right)
\end{equation}

$C^+$ characteristic slope:
\begin{equation}
    m^+ = \tan\!\left(\frac{\theta_{\mathrm{int}} + \theta_{\mathrm{wall},j}}{2} + \frac{\mu_{\mathrm{int}} + \mu_{\mathrm{wall},j}}{2}\right)
\end{equation}

Intersect these two lines:
\begin{align}
    y - y_{\mathrm{prev}} &= m_{\mathrm{wall}}(x - x_{\mathrm{prev}}) \\
    y - y_{\mathrm{int}} &= m^+(x - x_{\mathrm{int}})
\end{align}

Solve for $(x_{\mathrm{wall},j},\; y_{\mathrm{wall},j})$.

\section{Summary of the Algorithm}

\subsection{Phase 1: Setup}
\begin{enumerate}
    \item Compute $\nu_e$ from $M_e = 6$ using the Prandtl--Meyer function.
    \item $\theta_{\max} = \nu_e / 2$.
    \item $\Delta\theta = \theta_{\max} / N$.
    \item At each wall source point $i$ ($i = 1, \ldots, N$): $\theta_i = i\,\Delta\theta$, $\nu_i = i\,\Delta\theta$, $K^-_i = 2i\,\Delta\theta$, $K^+ = 0$.
    \item Compute $M_i$ and $\mu_i$ at each point from $\nu_i$.
    \item Compute $(x_i, y_i)$ on the circular arc.
\end{enumerate}

\subsection{Phase 2: March to Centerline (First Column)}
\begin{enumerate}
    \item The first $C^-$ characteristic (from wall point 1) travels to the centerline.
    \item At the centerline: $\theta = 0$, $\nu = K^-_1 = 2\Delta\theta$.
    \item Physically locate this point by intersecting the $C^-$ line with $y = 0$.
\end{enumerate}

\subsection{Phase 3: Interior Grid (Non-Simple Region)}
\begin{enumerate}
    \item Systematically compute all intersection points.
    \item At each point, use $K^-$ and $K^+$ from the two crossing characteristics.
    \item Compute $\theta$, $\nu$, $M$, $\mu$, and $(x, y)$.
\end{enumerate}

\subsection{Phase 4: Wall Contour (Straightening Section)}
\begin{enumerate}
    \item When $C^+$ characteristics reach the upper wall, apply $K^- = \nu_e$.
    \item Compute $\theta_{\mathrm{wall}}$ and $\nu_{\mathrm{wall}}$.
    \item Find $(x, y)$ by intersecting the $C^+$ line with the wall line from the previous point.
\end{enumerate}

\subsection{Phase 5: Verification}
At the exit plane, all points should have $\theta = 0$ and $\nu = \nu_e$ (uniform $M = 6$ flow, horizontal).

\section{Key Equations Reference}

\textbf{Prandtl--Meyer function} (for $\gamma = 1.4$):
\begin{equation}
    \nu(M) = \sqrt{6}\,\arctan\!\left(\sqrt{\frac{M^2 - 1}{6}}\right) - \arctan\!\left(\sqrt{M^2 - 1}\right)
\end{equation}

\textbf{General form:}
\begin{equation}
    \nu(M) = \sqrt{\frac{\gamma+1}{\gamma-1}}\,\arctan\!\left(\sqrt{\frac{(\gamma-1)(M^2-1)}{\gamma+1}}\right) - \arctan\!\left(\sqrt{M^2-1}\right)
\end{equation}

\textbf{Mach angle:}
\begin{equation}
    \mu = \arcsin\!\left(\frac{1}{M}\right)
\end{equation}

\textbf{Isentropic area ratio:}
\begin{equation}
    \frac{A}{A^*} = \frac{h}{h^*} = \frac{1}{M}\left(\frac{2}{\gamma+1}\left(1 + \frac{\gamma-1}{2}M^2\right)\right)^{\!\frac{\gamma+1}{2(\gamma-1)}}
\end{equation}

\end{document}